\section{Experimental setup}\label{sec:setup}

\subsection{Domain 1: graph path planning (exact validator)}
Each instance is a directed graph with 5, 10, 15 or 20 nodes, given in the prompt as one adjacency line per node,
together with a start and a goal node and one worked example. The agent returns the path as a JSON list. The
validator checks that the path starts and ends correctly and that every step $u\to v$ is an edge, i.e.\ that $v$
appears in the adjacency line of $u$. This validator is exact and costs microseconds, so this domain measures the
quality of the screens; compute savings are only meaningful in the second domain. Graphs are disjoint across all
datasets. Probes are fitted on a development set of 3000 graphs (1500 training, 500 calibration, 500 threshold
selection, 500 test) and evaluated once on a fresh certification set of 3000 graphs that had not been used before
(Table~\ref{tab:datasets}).

\subsection{Domain 2: CSTR fault recovery (digital-twin validator)}
The plant is the CSTR case study of \citep{car_todo}: a jacketed reactor with level, temperature and inlet-flow PID
loops, a monitoring layer that triggers the agent when the process leaves its safe operating region, and a digital
twin. Each episode draws a nominal plant (heat-transfer coefficient, feed and coolant temperatures, inlet
concentrations, coolant capacity, nominal level and feed setpoints), a fault family (jacket fouling, outlet pump
degradation, cooling valve stuck partly closed, outlet blockage), a continuous fault severity and an onset time
(1500--5000\,s). The plant is simulated to the first monitor trigger, and the snapshot at that moment is the agent's
task. Episodes that never trigger, or trigger before the fault onset, are redrawn within the same fault family.

The prompt contains the agent's role, the plant's knowledge graph in Turtle (process operators, controllers and the
mass and energy balance equations), the snapshot (measurements, actuator commands, alarms and current setpoints),
the goal and the admissible setpoint ranges. The prompt is deliberately \emph{neutral}: it states plant physics and
the acceptance criteria but gives no advice on which setpoint to change, in which direction or by how much, so that
failures reflect the model's own reasoning rather than prompt wording. The agent returns three setpoints (temperature,
level, feed flow) as JSON after a short justification.

The validator simulates the plant forward from the snapshot under the proposed setpoints, with the plant's own future
disturbances, to the end of the horizon. A proposal passes if the plant reaches the safe state (temperature at most
313\,K, level within 2--11, no overflow, no actuator saturated while the temperature is away from 310\,K) within
600\,s, holds it for 60 consecutive seconds, and remains safe for at least 80\% of the remaining run. One validation
took 7.8--11.6\,s on our hardware, depending on concurrent load.

Fault severities were calibrated without the LLM so that the reference model's (Qwen2.5-3B) first proposals pass in
roughly 30\% of episodes, which gives enough passing proposals to train and certify an accept decision. This is a
property of the dataset, not a target for the agents; pass rates of the other models range from 17\% to 35\%. The
5000 episodes are partitioned by episode into training (3000), calibration (500), threshold selection (500), test
(500) and certification (500), each with equal numbers per fault family. A further 400 fresh episodes, disjoint from
all others, are used for the closed loop.

\subsection{Models and inference}
Qwen2.5-1.5B, Qwen2.5-3B and Qwen2.5-7B \citep{qwen2025}, Llama-3.2-3B \citep{llama3} and SmolLM2-1.7B
\citep{allal2025smollm2} instruction-tuned models were run locally in bfloat16, except Qwen2.5-7B, which was run with
4-bit NF4 weights \citep{dettmers2023qlora} to fit an 8\,GB laptop GPU (NVIDIA RTX 4060). Decoding is greedy, so every
proposal is reproducible. The knowledge graph was exported once and frozen, so all prompts are identical across runs.

\subsection{Pre-registration and data hygiene}
For each study we froze a pre-registration before computing the relevant features or outcomes, recording later
changes as dated amendments made before the affected data were examined: the change of primary endpoint to validator
calls skipped, the bound-based accept rule, additional signal combinations, the closed-loop policies, the
first-proposal-only rule and a pilot gate for SmolLM2 in the CSTR. Test and certification partitions were evaluated
once per model with frozen probes and thresholds. Two analyses are labelled exploratory: the CSTR results of
Qwen2.5-3B under the bound-based rule (its sealed sets had been used once under the earlier rule), and the CSTR
closed-loop probes of Qwen2.5-3B, which come from that exploratory freeze (the closed-loop episodes themselves are
fresh). Where a run failed for technical reasons (GPU memory exhaustion, a crashed simulator worker), it was resumed
from its saved records; no records were discarded.

\begin{table}[t]
\centering
\small
\caption{Evaluation sets used for the confirmatory results (one candidate per graph or episode). FSM: fresh certification set, never used before evaluation. CSTR: sealed test and certification partitions, each evaluated once.}
\label{tab:datasets}
\begin{tabular}{llrrr}
\toprule
Domain & Model & Set & Eligible & Failure rate \\
\midrule
FSM & Qwen-3B & cert2 (fresh) & 2992 & 0.51 \\
FSM & Llama-3B & cert2 (fresh) & 2995 & 0.63 \\
FSM & Qwen-1.5B & cert2 (fresh) & 2945 & 0.68 \\
FSM & SmolLM-1.7B & cert2 (fresh) & 2946 & 0.83 \\
FSM & Qwen-7B & cert2 (fresh) & 2973 & 0.49 \\
CSTR & Qwen-1.5B & test_iid & 494 & 0.67 \\
CSTR & Qwen-1.5B & cert & 492 & 0.66 \\
CSTR & Qwen-3B$^\dagger$ & test_iid & 493 & 0.71 \\
CSTR & Qwen-3B$^\dagger$ & cert & 498 & 0.67 \\
CSTR & Qwen-7B & test_iid & 499 & 0.79 \\
CSTR & Qwen-7B & cert & 498 & 0.80 \\
CSTR & Llama-3B & test_iid & 473 & 0.84 \\
CSTR & Llama-3B & cert & 474 & 0.83 \\
CSTR & SmolLM-1.7B & test_iid & 499 & 0.67 \\
CSTR & SmolLM-1.7B & cert & 496 & 0.66 \\
\bottomrule
\end{tabular}
\end{table}

